\begin{filecontents}{refs.bib}
@article{Boerhout2023Microvascular,
title={Microvascular resistance reserve: diagnostic and prognostic performance in the ILIAS registry},
author={C. Boerhout and J. Lee and Guus A. de Waard and H. Mejía-Rentería and Seung Hun Lee and Ji-Hyun Jung and M. Hoshino and M. Echavarría-Pinto and M. Meuwissen and H. Matsuo and Maribel Madera-Cambero and A. Eftekhari and M. Effat and T. Murai and K. Marques and J. Doh and E. Christiansen and R. Banerjee and C. Nam and G. Niccoli and M. Nakayama and N. Tanaka and E. Shin and Y. Appelman and M. Beijk and N. van Royen and P. Knaapen and J. Escaned and T. Kakuta and B. Koo and J. Piek and T. P. van de Hoef},
journal={European Heart Journal},
year={2023},
volume={44},
pages={2862 - 2869},
doi={10.1093/eurheartj/ehad378}
}

@article{De Bruyne2021Microvascular,
title={Microvascular Resistance Reserve for Assessment of Coronary Microvascular Function: JACC Technology Corner.},
author={B. de Bruyne and N. Pijls and E. Gallinoro and A. Candreva and S. Fournier and D. Keulards and J. Sonck and M. Van't Veer and E. Barbato and J. Bartunek and M. Vanderheyden and E. Wyffels and A. D. de Vos and M. El Farissi and P. Tonino and O. Muller and C. Collet and W. Fearon},
journal={Journal of the American College of Cardiology},
year={2021},
volume={78 15},
pages={1541-1549},
doi={10.1016/j.jacc.2021.08.017}
}

@article{De Vos2023Microvascular,
title={Microvascular Resistance Reserve to Assess Microvascular Dysfunction in ANOCA Patients.},
author={A. D. de Vos and T. Jansen and M. van ’t Veer and A. Dimitriu-Leen and R. Konst and S. Elias-Smale and V. Paradies and L. Rodwell and S. V. D. van den Oord and P. Smits and Niels van Royen and N. Pijls and P. Damman},
journal={JACC. Cardiovascular interventions},
year={2023},
volume={16 4},
pages={470-481},
doi={10.1016/j.jcin.2022.12.012}
}

@article{Galante2025Fractional,
title={Fractional flow reserve (FFR) and index of microcirculatory resistance (IMR) relationship in patients with chronic or stabilized acute coronary syndromes.},
author={D. Galante and A. Viceré and A. Marrone and F. M. Verardi and V. Viccaro and C. Giuliana and C. P. Benvenuto and S. Todisco and S. Biscaglia and C. Aurigemma and Enrico Romagnoli and Gennaro Capalbo and C. Trani and F. Burzotta and Filippo Crea and Gianluca Campo and A. M. Leone},
journal={International journal of cardiology},
year={2025},
pages={132978},
doi={10.1016/j.ijcard.2025.132978}
}

@article{Garcia2019Relationship,
title={Relationship between FFR, CFR and coronary microvascular resistance – Practical implications for FFR-guided percutaneous coronary intervention},
author={Damien Garcia and B. Harbaoui and T. P. van de Hoef and M. Meuwissen and S. Nijjer and M. Echavarría-Pinto and J. Davies and J. Piek and P. Lantelme},
journal={PLoS ONE},
year={2019},
volume={14},
doi={10.1371/journal.pone.0208612}
}

@article{Jeremias2000Adequacy,
title={Adequacy of intracoronary versus intravenous adenosine-induced maximal coronary hyperemia for fractional flow reserve measurements.},
author={A. Jeremias and R. Whitbourn and S. Filardo and P. Fitzgerald and D. Cohen and E. Tuzcu and W. Anderson and A. Abizaid and G. Mintz and A. Yeung and M. Kern and P. Yock},
journal={American heart journal},
year={2000},
volume={140 4},
pages={651-7},
doi={10.1067/mhj.2000.109920}
}

@article{Johnson2015Repeatability,
title={Repeatability of Fractional Flow Reserve Despite Variations in Systemic and Coronary Hemodynamics.},
author={N. Johnson and Daniel T. Johnson and R. Kirkeeide and C. Berry and Mb Chb and B. de Bruyne and W. Fearon and K. Oldroyd and N. Pijls and K Lance Gould},
journal={JACC. Cardiovascular interventions},
year={2015},
volume={8 8},
pages={1018-1027},
doi={10.1016/j.jcin.2015.01.039}
}

@article{Johnson2016Continuum,
title={Continuum of Vasodilator Stress From Rest to Contrast Medium to Adenosine Hyperemia for Fractional Flow Reserve Assessment.},
author={MD MS Nils P. Johnson and MD Allen Jeremias and MD Frederik M. Zimmermann and MD Julien Adjedj and MD Nils Witt and MB Barry Hennigan and BC H Bao and CI Bmeds and MD Bon-Kwon Koo and MD Akiko Maehara and BS Mitsuaki Matsumura and MD Emanuele Barbato and MD Giovanni Esposito and MD Bruno Trimarco and MD Gilles Rioufol and MD Seung-Jung Park and MD Hyoung-Mo Yang and MD Sérgio B. Baptista and MD George S. Chrysant and MD Antonio M. Leone and Mbc H B Colin Berry and MD Bernard De Bruyne and MD K. Lance Gould and P. H. D. Richard L. Kirkeeide and Mbc H B Keith G. Oldroyd and MD Nico H.J. Pijls and MD William F. Fearon},
journal={JACC. Cardiovascular interventions},
year={2016},
volume={9 8},
pages={757-767},
doi={10.1016/j.jcin.2015.12.273}
}

@article{Leone2012Maximal,
title={Maximal hyperemia in the assessment of fractional flow reserve: intracoronary adenosine versus intracoronary sodium nitroprusside versus intravenous adenosine: the NASCI (Nitroprussiato versus Adenosina nelle Stenosi Coronariche Intermedie) study.},
author={A. Leone and I. Porto and A. D. De Caterina and Eloisa Basile and Andrea Aurelio and A. Gardi and D. Russo and D. Laezza and G. Niccoli and F. Burzotta and C. Trani and M. Mazzari and R. Mongiardo and A. Rebuzzi and F. Crea},
journal={JACC. Cardiovascular interventions},
year={2012},
volume={5 4},
pages={402-8},
doi={10.1016/j.jcin.2011.12.014}
}

@article{Morris2020A,
title={A novel method for measuring absolute coronary blood flow and microvascular resistance in patients with ischaemic heart disease},
author={Paul D Morris and R. Gosling and I. Zwierzak and Holli Evans and Louise Aubinière-Robb and K. Czechowicz and Paul C. Evans and D. Hose and Patricia V. Lawford and A. Narracott and Julian P. Gunn},
journal={Cardiovascular Research},
year={2020},
volume={117},
pages={1567 - 1577},
doi={10.1093/cvr/cvaa220}
}

@article{Soto2017Fast,
title={Fast Virtual Fractional Flow Reserve Based Upon Steady-State Computational Fluid Dynamics Analysis},
author={Daniel Alejandro Silva Soto and A. Lungu and Patricia V. Lawford},
journal={JACC: Basic to Translational Science},
year={2017},
volume={2},
pages={434 - 446},
doi={10.1016/j.jacbts.2017.04.003}
}

@article{Tu2019Fractional,
title={Fractional flow reserve in clinical practice: from wire-based invasive measurement to image-based computation.},
author={S. Tu and J. Westra and J. Adjedj and D. Ding and F. Liang and Bo Xu and N. Holm and J. Reiber and W. Wijns},
journal={European heart journal},
year={2019},
doi={10.1093/eurheartj/ehz918}
}

@article{Van De Hoef2014Impact,
title={Impact of hyperaemic microvascular resistance on fractional flow reserve measurements in patients with stable coronary artery disease: insights from combined stenosis and microvascular resistance assessment},
author={T. P. van de Hoef and F. Nolte and M. Echavarría-Pinto and M. V. van Lavieren and P. Damman and Steven A J Chamuleau and M. Voskuil and H. Verberne and J. Henriques and B. van Eck-Smit and K. Koch and R. D. de Winter and J. Spaan and M. Siebes and J. Tijssen and M. Meuwissen and J. Piek},
journal={Heart},
year={2014},
volume={100},
pages={951 - 959},
doi={10.1136/heartjnl-2013-305124}
}

@article{Verdoia2019Impact,
title={Impact of increasing dose of intracoronary adenosine on peak hyperemia duration during fractional flow reserve assessment.},
author={M. Verdoia and R. Erbetta and Emanuele Sagazio and L. Barbieri and F. Negro and H. Suryapranata and Elvin Kedhi and G. De Luca},
journal={International journal of cardiology},
year={2019},
volume={284},
pages={16-21},
doi={10.1016/j.ijcard.2018.09.088}
}
\end{filecontents}

\documentclass{article}
\usepackage{adjustbox}
\begin{document}

\subsection*{Key points from the literature}


The retrieved papers mainly address (1) why \textbf{maximal hyperemia with adenosine is required for FFR}, and (2) how \textbf{FFR depends on microvascular resistance} using both theoretical and in‑vivo data. None are the original first‑generation clinical validation trials (e.g., DEFER, FAME), but they clarify the physiological assumptions and limitations of FFR.


\subsubsection*{Why adenosine‑induced hyperemia is considered mandatory}


FFR is defined as the ratio of distal to aortic pressure during \textbf{maximal vasodilation}, so that coronary microvascular resistance is minimized and relatively stable; only under this condition does a pressure ratio reflect the relative reduction in \textbf{maximal myocardial flow} produced by a stenosis  \cite{Johnson2015Repeatability} \cite{Johnson2016Continuum}. Autoregulation at rest prevents pressure from linearly reflecting flow, hence the emphasis on hyperemia  \cite{Johnson2015Repeatability}.


Adenosine (intravenous or intracoronary) is the dominant agent because:


\begin{itemize}
\item It reliably produces \textbf{sustained, near‑maximal hyperemia} with a favorable safety profile, and became standard in pivotal trials such as FAME  \cite{Johnson2015Repeatability} \cite{Johnson2016Continuum}.
\item Consensus documents explicitly recommend \textbf{intravenous adenosine 140 μg/kg/min} as the reference method for clinical and trial use  \cite{Johnson2015Repeatability} \cite{Jeremias2000Adequacy} \cite{Johnson2016Continuum}.
\item Alternative vasodilators (e.g., intracoronary nitroprusside, papaverine, high‑dose intracoronary adenosine) can achieve comparable FFR values, but show more variability and, in the case of IC adenosine, a higher risk of \textbf{submaximal hyperemia} in a minority of patients  \cite{Leone2012Maximal} \cite{Jeremias2000Adequacy} \cite{Johnson2016Continuum} \cite{Verdoia2019Impact}.
\item Studies of the “hyperemic continuum” show that contrast‑induced \textbf{submaximal} hyperemia yields good, but clearly inferior, diagnostic performance vs. adenosine‑FFR and is proposed only as an adenosine‑sparing predictor, not a replacement; FFR with strong adenosine hyperemia remains the \textbf{reference standard}  \cite{Johnson2016Continuum}.
\end{itemize}


\subsubsection*{Sensitivity of FFR to microvascular resistance: models and in‑vivo data}


A combined theoretical and clinical analysis explicitly expressed FFR as a function of the \textbf{pressure‑loss coefficient of the stenosis, aortic pressure, and hyperemic coronary microvascular resistance (HMR)}, and showed:


\begin{itemize}
\item At intermediate severity, \textbf{HMR and stenosis pressure‑loss contributed comparably to FFR (≈45\% vs. 49\%)}  \cite{Garcia2019Relationship}.
\item FFR \textbf{increases with increasing HMR} for a given epicardial lesion (i.e., microvascular dysfunction can make a stenosis look less severe)  \cite{Garcia2019Relationship} \cite{Van De Hoef2014Impact}.
\item The mathematical FFR–CFR relationship is “exclusively dependent on coronary microvascular resistances,” and the model was validated in 199 lesions with simultaneous pressure/velocity measurements  \cite{Garcia2019Relationship}.
\end{itemize}


In‑vivo work that simultaneously measured FFR, hyperemic stenosis resistance (HSR) and HMR in 255 intermediate lesions confirmed these model predictions:


\begin{itemize}
\item The relationship between FFR and HSR (epicardial disease severity) was \textbf{strongly modulated by HMR}; after adjusting for HMR the FFR–HSR relation improved markedly  \cite{Van De Hoef2014Impact}.
\item For lesions with similar epicardial severity, \textbf{FFR was higher when HMR was higher}, demonstrating that FFR is partially “obscured” by downstream resistance  \cite{Van De Hoef2014Impact}.
\item The diagnostic odds ratio of FFR for ischemia on myocardial perfusion scintigraphy was significantly better when \textbf{HMR was high}, highlighting the interplay between stenosis and microcirculation  \cite{Van De Hoef2014Impact}.
\end{itemize}


Computational FFR work also confirms strong sensitivity to microvascular parameters: in a vFFR CFD framework, \textbf{microvascular physiology accounted for \textasciitilde{}59\% of the variance in physiological lesion significance}, vs. 33\% for coronary/lesion anatomy  \cite{Soto2017Fast}. A separate CFD‑based method that computes absolute flow and microvascular resistance alongside FFR explicitly emphasizes that pressure‑only indices cannot quantify microvascular disease and only predict relative, not absolute, flow change  \cite{Morris2020A}.


More recent registry data using microvascular resistance reserve (MRR) derive the index directly from CFR and FFR corrected for driving pressure, and show that as FFR falls, \textbf{discrepancies between CFR and microvascular‑specific indices widen}, again underscoring that FFR is not “microvascular‑independent”  \cite{De Bruyne2021Microvascular} \cite{Boerhout2023Microvascular} \cite{De Vos2023Microvascular}.


\subsubsection*{Is FFR independent of microvascular resistance in moderate stenosis?}


The convergent evidence from theory, CFD, and invasive measurements indicates that FFR is \textbf{not independent} of microvascular resistance, even in intermediate lesions:


\begin{itemize}
\item The explicit hemodynamic model and its validation show that FFR around the 0.80 threshold is jointly determined by stenosis severity and \textbf{hyperemic microvascular resistance}, with roughly equal impact  \cite{Garcia2019Relationship}.
\item In stable coronary disease, for stenoses of intermediate severity, \textbf{FFR systematically increased with higher HMR for the same epicardial lesion}, demonstrating dependence on Rs downstream  \cite{Van De Hoef2014Impact}.
\item Clinical registries demonstrate that elevated IMR/HMR alters FFR behavior: high microvascular resistance attenuates the hyperemic flow response and raises FFR, whereas very low microvascular resistance in large perfusion territories may drive \textbf{unusually low FFR} for only moderate epicardial disease  \cite{Tu2019Fractional} \cite{Galante2025Fractional} \cite{Van De Hoef2014Impact}.
\end{itemize}


None of these studies support strict independence of FFR from microvascular resistance in moderate stenosis; instead, they show a \textbf{bidirectional coupling}, with microvascular resistance materially influencing the measured FFR value.


\subsubsection*{Illustrative comparison of modeling vs in‑vivo evidence}


\mbox{}\\
\begin{adjustbox}{max width=\textwidth}
\begin{tabular}{p{4.0cm} | p{4.0cm} | p{4.0cm} | p{4.0cm}}
\hline
Aspect & Computational / theoretical & In‑vivo invasive data & Citations \\
\hline
FFR determinants & Explicit function of stenosis pressure‑loss, Pa, hyperemic microvascular resistance & HMR and stenosis resistance each explain large fractions of FFR variance & \cite{Soto2017Fast} \cite{Garcia2019Relationship} \cite{Van De Hoef2014Impact} \\
Contribution of microvasculature & \textasciitilde{}59\% of variance in vFFR significance attributed to microvascular physiology & FFR increases with higher HMR for same HSR; FFR–HSR relation improves after HMR adjustment & \cite{Soto2017Fast} \cite{Garcia2019Relationship} \cite{Van De Hoef2014Impact} \\
Independence from Rs? & No – equations make FFR explicitly Rs‑dependent & No – measured data confirm dependence around clinical thresholds & \cite{Garcia2019Relationship} \cite{Galante2025Fractional} \cite{Van De Hoef2014Impact} \\
\hline
\end{tabular}
\end{adjustbox}


\textbf{Figure 1:} Modeling and invasive data both show FFR depends on microvascular resistance.

% These search results were found and analyzed using Consensus, an AI-powered search engine for research. Try it at https://consensus.app. © 2026 Consensus NLP, Inc. Personal, non-commercial use only; redistribution requires copyright holders’ consent.

\bibliographystyle{plain}
\bibliography{refs}
\end{document}